% Autor: Elias Welt
\section{Deployment und Betrieb}
Für den produktiven Einsatz wird die Frontend-Anwendung containerisiert. Dies gewährleistet eine konsistente Laufzeitumgebung unabhängig vom Host-System.

\subsection{Docker-Containerisierung}
Es wird ein "`Multi-Stage Build"' Prozess verwendet, um das finale Image so klein wie möglich zu halten.
\begin{enumerate}
    \item \textbf{Build-Stage}: Ein Node.js Image wird verwendet, um die Abhängigkeiten zu installieren und den Build-Prozess (\texttt{npm run build}) auszuführen. Das Ergebnis ist ein Ordner \texttt{dist}, der nur die statischen Dateien (HTML, CSS, JS) enthält.
    \item \textbf{Production-Stage}: Ein leichtgewichtiger Nginx-Webserver wird als Basis genommen. Die Dateien aus der Build-Stage werden in das Verzeichnis \texttt{/usr/share/nginx/html} kopiert.
\end{enumerate}

\begin{lstlisting}[language=bash, caption={Dockerfile für das Frontend}]
# Stage 1: Build
FROM node:20-alpine as build
WORKDIR /app
COPY package*.json ./
RUN npm ci
COPY . .
RUN npm run build

# Stage 2: Serve
FROM nginx:alpine
COPY --from=build /app/dist /usr/share/nginx/html
COPY nginx.conf /etc/nginx/conf.d/default.conf
EXPOSE 80
CMD ["nginx", "-g", "daemon off;"]
\end{lstlisting}

\subsection{Nginx Konfiguration}
Da es sich um eine SPA handelt, muss der Webserver so konfiguriert werden, dass er für alle Routen (z.B. \texttt{/files}, \texttt{/chat}) immer die \texttt{index.html} ausliefert ("Fallback"). Andernfalls würde ein direkter Aufruf einer Unterseite zu einem 404-Fehler führen, da diese Dateien physisch nicht existieren.
